\section*{Supporting Information}

The Supporting Information is available free of charge. Summary tables are shown below; 
complete per-configuration data are provided as supplementary CSV files in the code 
repository.

\subsection*{Table S0. Endpoint classification criteria}

The full mapping of OECD, EPA, and EU Test Guidelines to the three endpoint categories 
(Ames, in~vitro CA, in~vivo MN) used for curating all datasets is provided as 
\texttt{table\_s0\_endpoint\_classification.xlsx}. This table was used to assign each raw 
study record to one of the target endpoints prior to any modeling.

\subsection*{Table S1. Random versus scaffold split comparison}

\begin{table}[ht]
\caption{Random vs.\ scaffold split (XGBoost, compact features, raw SMILES, 
$5\!\times\!5$ CV).}
\centering
\small
\begin{tabular}{@{}lrccc@{}}
\toprule
Endpoint    & $n$    & Random MCC        & Scaffold MCC      & $\Delta$ \\
\midrule
Ames        & 13,560 & $0.670 \pm 0.010$ & $0.624 \pm 0.049$ &  0.046   \\
In vitro CA &  1,619 & $0.180 \pm 0.053$ & $0.189 \pm 0.084$ & $-$0.009 \\
In vivo MN  &  2,153 & $0.181 \pm 0.094$ & $0.104 \pm 0.111$ &  0.077   \\
\bottomrule
\end{tabular}
\end{table}

\subsection*{Table S2. Full 225-experiment results}

Complete metrics (MCC, bootstrap CI, ROC-AUC, PR-AUC, balanced accuracy, sensitivity, 
specificity, Brier score, AD coverage) for all 225 configurations are provided in 
\texttt{table\_s2\_full\_results.csv}.

\subsection*{Table S3. Hyperparameter tuning comparison}

Default versus tuned MCC for all configurations, including optimized parameter values, are 
provided in \texttt{table\_s3\_hp\_tuning.csv}.

\subsection*{Table S4. Applicability domain coverage}

Per-configuration AD coverage (Tanimoto $\geq$ 0.3, 1024-bit Morgan FP) and mean similarity 
values are provided in \texttt{table\_s4\_ad\_coverage.csv}. Coverage ranged from 48.5\% 
(in~vitro) to 58.9\% (in~vivo) across configurations.

\subsection*{Table S5. Cross-endpoint compound overlap}

Pairwise compound overlap and label concordance ($\kappa$) between the three endpoint 
datasets are provided in \texttt{table\_s5\_cross\_endpoint.csv}. Overlap with Ames was 
97.5\% for in~vitro and 92.1\% for in~vivo.

\subsection*{Table S6. SHAP feature importance}

SHAP values (mean $|$SHAP$|$, TreeExplainer) for all features across all tree-based models 
are provided in \texttt{table\_s6\_shap.csv}.

\subsection*{Table S7. LDO threshold-tuning analysis (prior shift vs.\ source shift)}

To distinguish prior shift (calibration mismatch from differing prevalence) from genuine 
source shift, we evaluated each LDO direction at three decision thresholds: default 
($t = 0.5$), prior-corrected ($t$ = training positive rate), and oracle ($t$ maximizing MCC 
on the test set, providing an upper bound). Per-direction and per-algorithm results are 
shown below.

\textit{Note on features:} For computational tractability this analysis used Morgan FP-1024 
+ 13 physicochemical descriptors rather than the full 138-feature compact set used in the 
main pipeline. Absolute MCC values therefore differ from Table~\ref{tab:ldo} (notably RF 
industrial$\to$drug: 0.102 here vs.\ 0.325 in the main analysis). The internal 
three-threshold comparison within each row is valid because identical features are used 
across thresholds.

\begin{table}[ht]
\caption{LDO threshold-tuning analysis (Morgan FP-1024 + 13 physchem features). Default = 
MCC at $t = 0.5$. Prior-corrected = MCC at $t$ equal to the training positive rate. Oracle 
= upper bound MCC across all thresholds. The gap between oracle LDO MCC and the in-domain 
scaffold-CV MCC ($\approx$0.624) isolates the irreducible source shift that cannot be 
corrected by threshold tuning.}
\centering
\small
\begin{tabular}{@{}llcccc@{}}
\toprule
Direction & Algo. & MCC default & MCC prior-corr. & MCC oracle & Oracle threshold \\
\midrule
Drug $\to$ Ind.\ & XGBoost  & 0.231 & 0.263 & 0.310 & 0.788 \\
Drug $\to$ Ind.\ & LightGBM & 0.241 & 0.264 & 0.304 & 0.781 \\
Drug $\to$ Ind.\ & RF       & 0.250 & 0.299 & 0.359 & 0.690 \\
Ind.\ $\to$ Drug & XGBoost  & 0.200 & 0.267 & 0.273 & 0.041 \\
Ind.\ $\to$ Drug & LightGBM & 0.239 & 0.214 & 0.267 & 0.180 \\
Ind.\ $\to$ Drug & RF       & 0.102 & 0.341 & 0.345 & 0.037 \\
\bottomrule
\end{tabular}

\medskip
\footnotesize
Mean improvement (default $\to$ prior-corrected): $+0.075$ (range $-0.024$ to $+0.239$).
Mean improvement (default $\to$ oracle): $+0.108$ (range $+0.028$ to $+0.243$).
Oracle MCC range: 0.267--0.359, all substantially below the in-domain scaffold-CV MCC of 
0.624.
\end{table}
